


Generated by \texttt{python scripts/rd\_cv\_doc.py}. The checks of section 7 of \texttt{docs/02\_preregistration/PREREGISTRATION\_V12.md}; a module must pass them before it touches a fresh grid. Tests live in \texttt{tests/redesign/}.

\subsection*{1 to 3 and 7: tests (pass counts from the last run)}

{\scriptsize\raggedright
\begin{longtable}{p{0.190\linewidth}p{0.190\linewidth}p{0.190\linewidth}p{0.190\linewidth}p{0.190\linewidth}}
\toprule
Test file & Tests & Failed & Skipped & Covers \\
\midrule
\endhead
redesign & 39 & 0 & 0 &  \\
\textbf{all} & 39 & 0 & 0 &  \\
\bottomrule
\end{longtable}}

\subsection*{5: mutation testing (31 hand-written mutants; mutmut does not run on Windows)}

Mutants killed: \textbf{31 of 31} (100\%); the threshold is 80\%. Script \texttt{scripts/rd\_mutation.py}. The initial run killed 27 of 31 (87\%); the four survivors showed missing tests (table dependence on failed attempts, the crowd term, the cluster-ACK count check, the node-side ACK counter check), which were added, and the rerun is the number above.

{\scriptsize\raggedright
\begin{longtable}{p{0.237\linewidth}p{0.237\linewidth}p{0.237\linewidth}p{0.237\linewidth}}
\toprule
Id & File & Mutant & Result \\
\midrule
\endhead
M01 & rd\_policy.py & time left: sign of elapsed time & KILLED \\
M02 & rd\_policy.py & retry allowed: strict comparison & KILLED \\
M03 & rd\_policy.py & retry allowed: airtime ignored & KILLED \\
M04 & rd\_policy.py & Beta update: failure and success swapped & KILLED \\
M05 & rd\_policy.py & Beta forgetting skipped on a & KILLED \\
M06 & rd\_policy.py & Beta mean inverted & KILLED \\
M07 & rd\_policy.py & loss bin edge side & KILLED \\
M08 & rd\_policy.py & load bin edge side & KILLED \\
M09 & rd\_policy.py & budget units rounded up & KILLED \\
M10 & rd\_policy.py & lookup: no time left check removed & KILLED \\
M11 & rd\_policy.py & lookup ignores the number of failed attempts & KILLED \\
M12 & rd\_policy.py & DP feasibility ignores the gap & KILLED \\
M13 & rd\_policy.py & DP budget charge rounded down & KILLED \\
M14 & rd\_policy.py & DP crowd term removed & KILLED \\
M15 & rd\_policy.py & monotone projection over time removed & KILLED \\
M16 & rd\_policy.py & monotone projection over loss removed & KILLED \\
M17 & rd\_policy.py & cluster ACK length: one byte per id & KILLED \\
M18 & rd\_policy.py & signature not verified & KILLED \\
M19 & rd\_policy.py & event id check inverted & KILLED \\
M20 & rd\_policy.py & freshness window one-sided & KILLED \\
M21 & rd\_policy.py & replay check removed & KILLED \\
M22 & rd\_policy.py & processed frames not remembered & KILLED \\
M23 & rd\_policy.py & listing test inverted & KILLED \\
M24 & rd\_policy.py & length consistency check removed & KILLED \\
M25 & rd\_policy.py & vehicle does not stop on a verified listing & KILLED \\
M26 & rd\_policy.py & a defer becomes a stop & KILLED \\
M27 & rd\_policy.py & defer interval unbounded above and may lower ready & KILLED \\
M28 & verifier.py & base verifier: tag check skipped & KILLED \\
M29 & verifier.py & base verifier: replay counter off by one & KILLED \\
M30 & verifier.py & node: ACK counter not checked & KILLED \\
M31 & packet.py & tag comparison always true & KILLED \\
\bottomrule
\end{longtable}}

\subsection*{4: differential test (simulator policy against the firmware-path implementation)}

NOT\_RUN: the firmware-path implementation (D6) is built in session 2; the test is run then.

\subsection*{6: independent review pass}

A reviewer in a fresh context (a subagent told not to read results, tables or DEV logs) read the policy, simulator hooks, evaluation code and bounded model against \texttt{docs/02\_preregistration/PREREGISTRATION\_V12.md} before the freeze. It reported no blocker and no use of held-out data in the evaluation path (the old held-out logs are read only by \texttt{rd\_d0.py}, the declared DIAGNOSIS step). 20 findings and what was done:

{\scriptsize\raggedright
\begin{longtable}{p{0.237\linewidth}p{0.237\linewidth}p{0.237\linewidth}p{0.237\linewidth}}
\toprule
\# & Finding (short) & Severity & Disposition \\
\midrule
\endhead
1 & a demoted D1 alert stopped after one more try instead of the class limit (8 or 12) & MAJOR & fixed in \texttt{world\_sim.py} \\
2 & a simulated round is flash plus alert, the table plans one alert per attempt & MAJOR & disclosed (V12 addendum A, 5a); not changed \\
3 & event bootstrap ignores correlation within a simulated hour & MAJOR & verdicts now use the more conservative of the event and the hour-cluster bootstrap (addendum A, 1) \\
4 & H-D2c model run with smaller bounds than preregistered & MAJOR & 2 vehicles now at the preregistered bounds; 3 vehicles at reduced bounds, stated (addendum A, 3) \\
5 & the stolen-gateway-key limit case could never be flagged & MAJOR & fixed in \texttt{rd\_d2\_model.py} \\
6 & two model properties hold by construction; the attacker flood needs per-vehicle keys & MINOR & disclosed in the model docstring and addendum A, 3 \\
7 & criterion form (p below 0.05 and point at least 0.05) differs from the two-sided interval & MINOR & stated in addendum A, 2 \\
8 & a defer could push a retry past the class deadline & MINOR & fixed (capped at deadline minus one round) \\
9 & the Beta mean and the table both count the current failures & MINOR & disclosed (5b) \\
10 & table caps NB\_MAX and J\_MAX have no stated origin and bind & MINOR & disclosed (5c) \\
11 & the cluster ACK lists only the 24 most recent vehicles & MINOR & disclosed (5d) \\
12 & a demoted alert kept the raised SF & MINOR & fixed \\
13 & float32 delays at the deadline boundary & MINOR & fixed (float64) \\
14 & gate computed before the current frame is counted; the initial attempt uses only the table SF & NOTE & noted; unchanged \\
15 & Beta success counted once per alert, failures per round & NOTE & noted; unchanged \\
16 & the simulator does not call \texttt{ClusterAckVerifier} or \texttt{NodeState} & NOTE & covered by the D6 differential test (session 2) (5e) \\
17 & b = 1 cells contain a one-vehicle burst cell; D2 equals D1 there only because thr is at least 2 & NOTE & assertion added (thr of at least 2) \\
18 & constants without a stated origin & NOTE & listed (5g) and in \texttt{ASSUMPTIONS.md} \\
19 & DEV and FRESH-A share the finite class pool & NOTE & disclosed (5f) \\
20 & freeze hashed too few modules & NOTE & extended (addendum A, 6) \\
\bottomrule
\end{longtable}}

\subsection*{Findings of the checks so far}

\begin{itemize}
\item The property test of table monotonicity found 107 violations in the raw dynamic-programme output; a monotone projection was added (\texttt{D1Table.\_monotone}) and the test passes.
\item The invariant check found a variable shadowing bug in the initial version of the D2 hook (\texttt{active} hid the list of transmissions) and, in the initial D2 test, a cluster mode that triggered on one vehicle's repeated frames; the hook now counts distinct vehicles.
\item The invariant \texttt{events conserved} needed the event that ended the run to be kept in the queue so that transmissions in flight are counted.
\end{itemize}


\subsection*{9. Session 2: D3, D4, the C firmware path, the ns-3 application and FRESH-B machinery}

Tests: \texttt{tests/redesign/} now has 73 passing tests (session 1: 39). New files: \texttt{test\_rd\_d3.py} (hand-computed slot cases, property tests with hypothesis for the slot bounds, deadline and no-stop rule, simulator invariants, D3-off equals D1 + D2 bit for bit, simulator decisions equal the pure function), \texttt{test\_rd\_firmware\_diff.py} (SHA-256 and HMAC known answers including RFC 4231 case 3, airtime and bin hand cases, non-canonical Ed25519 signature rejected by both sides, a 1,000-trace differential run, replay of simulator decisions), \texttt{test\_rd\_d4\_fresh\_b.py} (hand-computed extra airtime, \textbf{known-answer reproduction of the stored R3 extra airtime 0.38038 and 0.33617}, monotonicity in the provisional frame length, a smoke run of the whole FRESH-B analysis on a tiny grid before the single real run), \texttt{test\_rd\_regression\_a.py} (\textbf{D1 and D1 + D2 under the current simulator reproduce the stored FRESH-A per-event delays bit for bit}), \texttt{test\_rd\_bench.py} (bench-log parser on a synthetic fixture, DUMMY DATA).

\textbf{CV4 differential test.} 10000 random vehicle histories (409999 operations: Beta updates, D1 lookups, bins, defers, D3 slots, the D2 decision, alert frames and tags, tampered and extended tags, freshness sequences, cluster ACKs of 16 kinds including valid, tampered, replayed, stale, truncated, malformed but signed, wrong key and non-canonical signature) compared between the C path and the Python reference: \textbf{0 mismatches}; 5140 decisions recorded inside real simulator runs replayed through the C path: \textbf{0 mismatches}. Weak spots found by the review and not fixed: some comparisons are weak by construction (the D2 decision is compared with a copy of the branch condition; the cluster check of the simulator is an abstraction of the signature check, so cryptography is tested by unit tests and the differential run, not inside the simulator); D3 slot times in the random ops use t\_min below 40 s (real times come from the replay only); the replay contains few D3 slot decisions (the D3 slots are rarely used, which is also the finding of H-D3).

\textbf{CV5 mutation testing} (\texttt{scripts/rd\_mutation\_b.py}, 22 mutants of \texttt{rd\_d3.py}, 5 of its simulator integration, 26 of the C path; mutmut does not run on Windows): 49 of 53 killed (92\%); survivors: D18,W04,C12,C26. A mutant counts as killed only when it compiles and a test or a differential comparison fails (an earlier run of this script counted compile errors caused by a space in the temporary path as kills; that was found and corrected, and the table was re-made). Survivors read one by one (\texttt{results/tables/rd\_mutation\_b.csv}): W04 (a vehicle "hears itself": \texttt{node\_link} returns false for the same node) and C26 (a no-op) are equivalent mutants; D18 (the 5 ms clock offset ignored when choosing the epoch) is equivalent except at an epoch boundary that no test hits; \textbf{C12 is a real weakness}: without the "no time left" check the C lookup reads outside the table for a non-positive time left, and the random traces happened to return the same "no action" answer; a test with a non-positive time left that checks the index is in range is not written.

\textbf{CV6 independent review} (fresh-context subagent, told not to read results; it read V12 sections 1, 3, 4, 6, 9, 10, the D3, D4, FRESH-B and firmware code and the tests). 18 findings, no use of held-out data, no way to run FRESH-B twice. Fixed before the freeze: H-D4b wording in the analysis, an analyse-only entry point and a pre-flight check for the single run, the freeze fallback when no D3 iteration meets the DEV floor and a longer freeze hash list, C hardening (exact tag frame length, ring bounds, harness token handling, HMAC over-length path), seeded Ed25519 keys in the generator, a stale-binary risk in a test. Disclosed and not changed because D3 DEV had already run (V12 10.6): the D3 gate counts other vehicles; the overhearing model; the slot is applied to the flash of a round; the 5 s epoch cap does not hold at SF12; S\_MAX; \texttt{alt} outside the budget plan; \texttt{world\_sim.py} and \texttt{rd\_eval.py} changed after the FRESH-A freeze (equivalence tested bit for bit). After the freeze nothing was changed.

\textbf{CV7 known answers.} The R3 extra airtime (above) and the FRESH-A D1 and D1 + D2 delays are reproduced with the current code. The three session 1 known answers (UA-SCT cell, OS-DR cell, R3) are in \texttt{test\_rd\_sim.py} and pass.

\textbf{Not done:} a differential test of the ns-3 application against the Python simulator (they are different simulators by design; the D1 lookup inside ns-3 is the same C code, tested above); mutation testing of \texttt{rd\_x13.py}; ESP32 or emulator execution of the C path.

